% Options for packages loaded elsewhere
% Options for packages loaded elsewhere
\PassOptionsToPackage{unicode}{hyperref}
\PassOptionsToPackage{hyphens}{url}
%
\documentclass[
  14pt,
  ignorenonframetext,
  aspectratio=169,
]{beamer}
\newif\ifbibliography
\usepackage{pgfpages}
\setbeamertemplate{caption}[numbered]
\setbeamertemplate{caption label separator}{: }
\setbeamercolor{caption name}{fg=normal text.fg}
\beamertemplatenavigationsymbolsempty
% remove section numbering
\setbeamertemplate{part page}{
  \centering
  \begin{beamercolorbox}[sep=16pt,center]{part title}
    \usebeamerfont{part title}\insertpart\par
  \end{beamercolorbox}
}
\setbeamertemplate{section page}{
  \centering
  \begin{beamercolorbox}[sep=12pt,center]{section title}
    \usebeamerfont{section title}\insertsection\par
  \end{beamercolorbox}
}
\setbeamertemplate{subsection page}{
  \centering
  \begin{beamercolorbox}[sep=8pt,center]{subsection title}
    \usebeamerfont{subsection title}\insertsubsection\par
  \end{beamercolorbox}
}
% Prevent slide breaks in the middle of a paragraph
\widowpenalties 1 10000
\raggedbottom
\usepackage{iftex}
\ifPDFTeX
  \usepackage[T1]{fontenc}
  \usepackage[utf8]{inputenc}
  \usepackage{textcomp} % provide euro and other symbols
\else % if luatex or xetex
  \usepackage{unicode-math} % this also loads fontspec
  \defaultfontfeatures{Scale=MatchLowercase}
  \defaultfontfeatures[\rmfamily]{Ligatures=TeX,Scale=1}
\fi
\usepackage{lmodern}

\usetheme[]{monash}
\ifPDFTeX\else
  % xetex/luatex font selection
\fi
% Use upquote if available, for straight quotes in verbatim environments
\IfFileExists{upquote.sty}{\usepackage{upquote}}{}
\IfFileExists{microtype.sty}{% use microtype if available
  \usepackage[]{microtype}
  \UseMicrotypeSet[protrusion]{basicmath} % disable protrusion for tt fonts
}{}
\makeatletter
\@ifundefined{KOMAClassName}{% if non-KOMA class
  \IfFileExists{parskip.sty}{%
    \usepackage{parskip}
  }{% else
    \setlength{\parindent}{0pt}
    \setlength{\parskip}{6pt plus 2pt minus 1pt}}
}{% if KOMA class
  \KOMAoptions{parskip=half}}
\makeatother

\usepackage{color}
\usepackage{fancyvrb}
\newcommand{\VerbBar}{|}
\newcommand{\VERB}{\Verb[commandchars=\\\{\}]}
\DefineVerbatimEnvironment{Highlighting}{Verbatim}{commandchars=\\\{\}}
% Add ',fontsize=\small' for more characters per line
\usepackage{framed}
\definecolor{shadecolor}{RGB}{248,248,248}
\newenvironment{Shaded}{\begin{snugshade}}{\end{snugshade}}
\newcommand{\AlertTok}[1]{\textcolor[rgb]{0.94,0.16,0.16}{#1}}
\newcommand{\AnnotationTok}[1]{\textcolor[rgb]{0.56,0.35,0.01}{\textbf{\textit{#1}}}}
\newcommand{\AttributeTok}[1]{\textcolor[rgb]{0.13,0.29,0.53}{#1}}
\newcommand{\BaseNTok}[1]{\textcolor[rgb]{0.00,0.00,0.81}{#1}}
\newcommand{\BuiltInTok}[1]{#1}
\newcommand{\CharTok}[1]{\textcolor[rgb]{0.31,0.60,0.02}{#1}}
\newcommand{\CommentTok}[1]{\textcolor[rgb]{0.56,0.35,0.01}{\textit{#1}}}
\newcommand{\CommentVarTok}[1]{\textcolor[rgb]{0.56,0.35,0.01}{\textbf{\textit{#1}}}}
\newcommand{\ConstantTok}[1]{\textcolor[rgb]{0.56,0.35,0.01}{#1}}
\newcommand{\ControlFlowTok}[1]{\textcolor[rgb]{0.13,0.29,0.53}{\textbf{#1}}}
\newcommand{\DataTypeTok}[1]{\textcolor[rgb]{0.13,0.29,0.53}{#1}}
\newcommand{\DecValTok}[1]{\textcolor[rgb]{0.00,0.00,0.81}{#1}}
\newcommand{\DocumentationTok}[1]{\textcolor[rgb]{0.56,0.35,0.01}{\textbf{\textit{#1}}}}
\newcommand{\ErrorTok}[1]{\textcolor[rgb]{0.64,0.00,0.00}{\textbf{#1}}}
\newcommand{\ExtensionTok}[1]{#1}
\newcommand{\FloatTok}[1]{\textcolor[rgb]{0.00,0.00,0.81}{#1}}
\newcommand{\FunctionTok}[1]{\textcolor[rgb]{0.13,0.29,0.53}{\textbf{#1}}}
\newcommand{\ImportTok}[1]{#1}
\newcommand{\InformationTok}[1]{\textcolor[rgb]{0.56,0.35,0.01}{\textbf{\textit{#1}}}}
\newcommand{\KeywordTok}[1]{\textcolor[rgb]{0.13,0.29,0.53}{\textbf{#1}}}
\newcommand{\NormalTok}[1]{#1}
\newcommand{\OperatorTok}[1]{\textcolor[rgb]{0.81,0.36,0.00}{\textbf{#1}}}
\newcommand{\OtherTok}[1]{\textcolor[rgb]{0.56,0.35,0.01}{#1}}
\newcommand{\PreprocessorTok}[1]{\textcolor[rgb]{0.56,0.35,0.01}{\textit{#1}}}
\newcommand{\RegionMarkerTok}[1]{#1}
\newcommand{\SpecialCharTok}[1]{\textcolor[rgb]{0.81,0.36,0.00}{\textbf{#1}}}
\newcommand{\SpecialStringTok}[1]{\textcolor[rgb]{0.31,0.60,0.02}{#1}}
\newcommand{\StringTok}[1]{\textcolor[rgb]{0.31,0.60,0.02}{#1}}
\newcommand{\VariableTok}[1]{\textcolor[rgb]{0.00,0.00,0.00}{#1}}
\newcommand{\VerbatimStringTok}[1]{\textcolor[rgb]{0.31,0.60,0.02}{#1}}
\newcommand{\WarningTok}[1]{\textcolor[rgb]{0.56,0.35,0.01}{\textbf{\textit{#1}}}}

\usepackage{longtable,booktabs,array}
\newcounter{none} % for unnumbered tables
\usepackage{calc} % for calculating minipage widths
\usepackage{caption}
% Make caption package work with longtable
\makeatletter
\def\fnum@table{\tablename~\thetable}
\makeatother
\usepackage{graphicx}
\makeatletter
\newsavebox\pandoc@box
\newcommand*\pandocbounded[1]{% scales image to fit in text height/width
  \sbox\pandoc@box{#1}%
  \Gscale@div\@tempa{\textheight}{\dimexpr\ht\pandoc@box+\dp\pandoc@box\relax}%
  \Gscale@div\@tempb{\linewidth}{\wd\pandoc@box}%
  \ifdim\@tempb\p@<\@tempa\p@\let\@tempa\@tempb\fi% select the smaller of both
  \ifdim\@tempa\p@<\p@\scalebox{\@tempa}{\usebox\pandoc@box}%
  \else\usebox{\pandoc@box}%
  \fi%
}
% Set default figure placement to htbp
\def\fps@figure{htbp}
\makeatother





\setlength{\emergencystretch}{3em} % prevent overfull lines

\providecommand{\tightlist}{%
  \setlength{\itemsep}{0pt}\setlength{\parskip}{0pt}}



 


% Colors
\definecolor{shadecolor}{RGB}{225,225,225}
\setbeamercolor{description item}{fg=Orange}
\definecolor{MonashBlue}{RGB}{66,109,152}
\definecolor{burntorange}{rgb}{0.8, 0.33, 0.0}

% Packages
\usepackage{amsmath, bm, amssymb, amsthm, mathrsfs,pifont}
\usepackage{url}
\usepackage{multirow, booktabs, float, textcmds, siunitx}
\usepackage{bm,booktabs,animate,ragged2e,multicol,microtype,hyperref}
\usepackage{fontawesome5}

% Figures
\graphicspath{{figs/}}

% Fonts
\fontsize{13}{15}\sf
\setbeamerfont{title}{series=\bfseries,parent=structure,size=\fontsize{24}{24}}
\ifcsname Shaded\endcsname
  \definecolor{shadecolor}{RGB}{225,225,225}
  \renewenvironment{Shaded}{\vspace*{0.15cm}\color{black}\fontsize{10}{10}\sf\begin{snugshade}\color{black}}{\end{snugshade}}
\fi

% gt dependencies
\usepackage{longtable,caption,setspace}
\captionsetup{font={small,stretch=0.80}}

% My definitions

\def\E{\text{E}}
\def\V{\text{Var}}
\def\up#1{\raisebox{-0.3cm}{#1}}
\def\pred#1#2#3{\hat{#1}_{#2|#3}}
\def\damped{$_\text{d}$}
\def\h+{h_{m}^{+}}
\def\str#1{\rlap{#1}\textcolor{red}{\rule{1cm}{0.1cm}}}

\def\st#1{\rlap{#1}\textcolor{red}{\rule{1cm}{0.1cm}}}
\def\bY{\bm{y}}
\def\by{\bm{y}}
\def\bS{\bm{S}}
\def\bI{\text{\rm\textbf{I}}}
\def\bbeta{\bm{\beta}}
\def\bSigma{\bm{\Sigma}}
\def\bW{\bm{\Sigma}}
\def\Var{\text{Var}}
\def\var{\text{Var}}
\def\bOmega{\bm{\Omega}}
\def\bLambda{\bm{\Lambda}}
\let\mc\multicolumn
\def\hl{\color[RGB]{230, 172, 0}}

\def\forecast{\begin{alertblock}{}\fontsize{10}{11}\sf
 A forecast is an estimate of the probability distribution of a variable to be observed in the future.
\end{alertblock}}

\def\simfutures{\begin{textblock}{2.9}(12.8,7.7)
\begin{block}{}\fontsize{10}{11}\sf
Simulated futures from an ETS model
\end{block}\end{textblock}}

\setbeamertemplate{title page}
{\placefig{0}{-0.5}{width=1.01\paperwidth,height=1.45\paperheight}{africa1.jpg}
\begin{textblock}{8.5}(.5,0.2)
  \raggedright\fontsize{20}{20}\selectfont\bfseries\sffamily\textcolor{white}{\inserttitle}\\[0.2cm]
  \fontsize{12}{12}\selectfont\bfseries\sffamily\textcolor{white}{\insertsubtitle}
  \end{textblock}
\placefig{.5}{6.2}{width=2.4cm}{tsibble.png}
\placefig{2.9}{6.2}{width=2.4cm}{feasts.png}
\placefig{5.3}{6.2}{width=2.4cm}{fable.png}
\begin{textblock}{7.5}(10.6,-.1)
  {\fontsize{9}{9}\sf\color{white}https://workshop.f4sg.org/africast/}
\end{textblock}
}

% Tikz plots
\usepackage{tikz}
\usetikzlibrary{trees,shapes,arrows,matrix,shadows,positioning,calc}
\tikzstyle{decision} = [diamond, draw, fill=blue!20,
    text width=4.5em, text badly centered, node distance=4cm, inner sep=0pt]
\tikzstyle{block} = [rectangle, draw, fill=blue!20,
    text width=5cm, text centered, rounded corners, minimum height=4em]
\tikzstyle{line} = [draw, thick, -latex']
\tikzstyle{cloud} = [draw, ellipse,fill=red!20, node distance=3cm,
    minimum height=2em, text centered]
\tikzstyle{connector} = [->,thick]

\tikzset{
  basic/.style  = {draw, text width=2cm, font=\sffamily, rectangle},
  root/.style   = {basic, text width=3cm, rounded corners=2pt, thin, align=center, fill=red!30},
  level 2a/.style = {basic, rounded corners=2pt, thin,align=center, fill=blue!50, text width=7em},
  level 2b/.style = {basic, rounded corners=2pt, thin,align=center, fill=green!50, text width=7em},
  level 3a/.style = {basic, rounded corners=2pt, thin, align=center, fill=blue!30, text width=4em},
  level 3b/.style = {basic, rounded corners=2pt, thin, align=center, fill=green!30, text width=4em},
  level 4a/.style = {basic, rounded corners=2pt, thin, align=left, fill=blue!10, text width=3.5em},
  level 4b/.style = {basic, rounded corners=2pt, thin, align=left, fill=green!10, text width=3.5em}
}

% % BIBLIOGRAPHIES
% \usepackage[style=authoryear,bibencoding=utf8,minnames=1,maxnames=4, maxbibnames=99,natbib=true,dashed=false,terseinits=true,giveninits=true,uniquename=false,uniquelist=false,labeldate=true,doi=false, isbn=false, natbib=true,backend=biber]{biblatex}

% \DeclareFieldFormat{url}{\url{#1}}
% \DeclareFieldFormat[article]{pages}{#1}
% \DeclareFieldFormat[inproceedings]{pages}{\lowercase{pp.}#1}
% \DeclareFieldFormat[incollection]{pages}{\lowercase{pp.}#1}
% \DeclareFieldFormat[article]{volume}{\mkbibbold{#1}}
% \DeclareFieldFormat[article]{number}{\mkbibparens{#1}}
% \DeclareFieldFormat[article]{title}{\MakeCapital{#1}}
% \DeclareFieldFormat[article]{url}{}
% \DeclareFieldFormat[Techreport]{Url}{}
% \DeclareFieldFormat[book]{url}{}
% \DeclareFieldFormat[inbook]{url}{}
% \DeclareFieldFormat[incollection]{url}{}
% \DeclareFieldFormat[inproceedings]{url}{}
% \DeclareFieldFormat[inproceedings]{title}{#1}
% \DeclareFieldFormat{shorthandwidth}{#1}
% %\DeclareFieldFormat{extrayear}{}
% % No dot before number of articles
% \usepackage{xpatch}
% \xpatchbibmacro{volume+number+eid}{\setunit*{\adddot}}{}{}{}
% % Remove In: for an article.
% \renewbibmacro{in:}{%
%   \ifentrytype{article}{}{%
%   \printtext{\bibstring{in}\intitlepunct}}}

% \AtEveryBibitem{\clearfield{month}}
% \AtEveryCitekey{\clearfield{month}}
% \AtBeginBibliography{\fontsize{11}{11}\sf}
\setbeamertemplate{frametitle continuation}{}
\makeatletter
\@ifpackageloaded{caption}{}{\usepackage{caption}}
\AtBeginDocument{%
\ifdefined\contentsname
  \renewcommand*\contentsname{Table of contents}
\else
  \newcommand\contentsname{Table of contents}
\fi
\ifdefined\listfigurename
  \renewcommand*\listfigurename{List of Figures}
\else
  \newcommand\listfigurename{List of Figures}
\fi
\ifdefined\listtablename
  \renewcommand*\listtablename{List of Tables}
\else
  \newcommand\listtablename{List of Tables}
\fi
\ifdefined\figurename
  \renewcommand*\figurename{Figure}
\else
  \newcommand\figurename{Figure}
\fi
\ifdefined\tablename
  \renewcommand*\tablename{Table}
\else
  \newcommand\tablename{Table}
\fi
}
\@ifpackageloaded{float}{}{\usepackage{float}}
\floatstyle{ruled}
\@ifundefined{c@chapter}{\newfloat{codelisting}{h}{lop}}{\newfloat{codelisting}{h}{lop}[chapter]}
\floatname{codelisting}{Listing}
\newcommand*\listoflistings{\listof{codelisting}{List of Listings}}
\makeatother
\makeatletter
\makeatother
\makeatletter
\@ifpackageloaded{caption}{}{\usepackage{caption}}
\@ifpackageloaded{subcaption}{}{\usepackage{subcaption}}
\makeatother

\usepackage{bookmark}
\IfFileExists{xurl.sty}{\usepackage{xurl}}{} % add URL line breaks if available
\urlstyle{same}
% fallback for those not using the hyperref driver hyperxmp:
\makeatletter
\@ifundefined{xmpquote}{\newcommand{\xmpquote}[1]{#1}}{}
\makeatother
\hypersetup{
  pdftitle={Africast-Time Series Analysis \& Forecasting Using R},
  hidelinks,
  pdfcreator={LaTeX via pandoc}}


\title{\texorpdfstring{Africast-Time Series Analysis \& Forecasting
Using R}{Africast-Time Series Analysis \& Forecasting Using R}}
\subtitle{\texorpdfstring{3. Computing and visualising
features}{3. Computing and visualising features}}
\author{}
\date{7 October 2026}

\begin{document}
\frame{\titlepage}


\begin{frame}{Outline}
\protect\phantomsection\label{outline}
\vspace*{0.7cm}\tableofcontents
\end{frame}

\section{STL Features}\label{stl-features}

\begin{frame}{Strength of seasonality and trend}
\protect\phantomsection\label{strength-of-seasonality-and-trend}
\begin{alertblock}{STL decomposition}
\centerline{$y_t = T_t+S_t+R_t$}
\end{alertblock}

\begin{block}{Seasonal strength}
\protect\phantomsection\label{seasonal-strength}
\[\max\left(0, 1-\frac{\text{Var}(R_t)}{\text{Var}(S_t+R_t)}\right)\]
\end{block}

\begin{block}{Trend strength}
\protect\phantomsection\label{trend-strength}
\[\max\left(0, 1-\frac{\text{Var}(R_t)}{\text{Var}(T_t+R_t)}\right)\]
\end{block}
\end{frame}

\begin{frame}[fragile]{Feature extraction and statistics}
\protect\phantomsection\label{feature-extraction-and-statistics}
\fontsize{9}{10}\sf

\begin{Shaded}
\begin{Highlighting}[]
\NormalTok{tourism }\SpecialCharTok{|\textgreater{}} \FunctionTok{features}\NormalTok{(Trips, feat\_stl)}
\end{Highlighting}
\end{Shaded}

\begin{verbatim}
# A tibble: 304 x 12
   Region         State Purpose trend_strength seasonal_strength_year seasonal_peak_year
   <chr>          <chr> <chr>            <dbl>                  <dbl>              <dbl>
 1 Adelaide       SA    Busine~          0.464                  0.407                  3
 2 Adelaide       SA    Holiday          0.554                  0.619                  1
 3 Adelaide       SA    Other            0.746                  0.202                  2
 4 Adelaide       SA    Visiti~          0.435                  0.452                  1
 5 Adelaide Hills SA    Busine~          0.464                  0.179                  3
 6 Adelaide Hills SA    Holiday          0.528                  0.296                  2
 7 Adelaide Hills SA    Other            0.593                  0.404                  2
 8 Adelaide Hills SA    Visiti~          0.488                  0.254                  0
 9 Alice Springs  NT    Busine~          0.534                  0.251                  0
10 Alice Springs  NT    Holiday          0.381                  0.832                  3
# i 294 more rows
# i 6 more variables: seasonal_trough_year <dbl>, spikiness <dbl>, linearity <dbl>,
#   curvature <dbl>, stl_e_acf1 <dbl>, stl_e_acf10 <dbl>
\end{verbatim}
\end{frame}

\begin{frame}[fragile]{Feature extraction and statistics}
\protect\phantomsection\label{feature-extraction-and-statistics-1}
\fontsize{8}{9}\sf

\begin{Shaded}
\begin{Highlighting}[]
\NormalTok{tourism }\SpecialCharTok{|\textgreater{}}
  \FunctionTok{features}\NormalTok{(Trips, feat\_stl) }\SpecialCharTok{|\textgreater{}}
  \FunctionTok{ggplot}\NormalTok{(}\FunctionTok{aes}\NormalTok{(}\AttributeTok{x =}\NormalTok{ trend\_strength, }\AttributeTok{y =}\NormalTok{ seasonal\_strength\_year, }\AttributeTok{col =}\NormalTok{ Purpose)) }\SpecialCharTok{+}
  \FunctionTok{geom\_point}\NormalTok{() }\SpecialCharTok{+} \FunctionTok{facet\_wrap}\NormalTok{(}\FunctionTok{vars}\NormalTok{(State))}
\end{Highlighting}
\end{Shaded}

\pandocbounded{\includegraphics[keepaspectratio]{02_features_files/figure-beamer/features-plot-1.pdf}}

\only<2->{\begin{textblock}{4.7}(9.8,6.9)
\begin{alertblock}{}\fontsize{10}{10}\sf
\begin{itemize}\tightlist
\item Holidays more seasonal than other travel.
\item WA has strongest trends.
\end{itemize}
\end{alertblock}\end{textblock}}
\end{frame}

\begin{frame}[fragile]{Feature extraction and statistics}
\protect\phantomsection\label{feature-extraction-and-statistics-2}
\fontsize{9}{9}\sf

Find the most seasonal time series:

\begin{Shaded}
\begin{Highlighting}[]
\NormalTok{most\_seasonal }\OtherTok{\textless{}{-}}\NormalTok{ tourism }\SpecialCharTok{|\textgreater{}}
  \FunctionTok{features}\NormalTok{(Trips, feat\_stl) }\SpecialCharTok{|\textgreater{}}
  \FunctionTok{filter}\NormalTok{(seasonal\_strength\_year }\SpecialCharTok{==} \FunctionTok{max}\NormalTok{(seasonal\_strength\_year))}
\end{Highlighting}
\end{Shaded}

\pause\vspace*{-0.3cm}

\begin{Shaded}
\begin{Highlighting}[]
\NormalTok{tourism }\SpecialCharTok{|\textgreater{}}
  \FunctionTok{right\_join}\NormalTok{(most\_seasonal, }\AttributeTok{by =} \FunctionTok{c}\NormalTok{(}\StringTok{"State"}\NormalTok{, }\StringTok{"Region"}\NormalTok{, }\StringTok{"Purpose"}\NormalTok{)) }\SpecialCharTok{|\textgreater{}}
  \FunctionTok{ggplot}\NormalTok{(}\FunctionTok{aes}\NormalTok{(}\AttributeTok{x =}\NormalTok{ Quarter, }\AttributeTok{y =}\NormalTok{ Trips)) }\SpecialCharTok{+}
  \FunctionTok{geom\_line}\NormalTok{() }\SpecialCharTok{+} \FunctionTok{facet\_grid}\NormalTok{(}\FunctionTok{vars}\NormalTok{(State, Region, Purpose))}
\end{Highlighting}
\end{Shaded}

\pandocbounded{\includegraphics[keepaspectratio]{02_features_files/figure-beamer/extreme2-1.pdf}}
\end{frame}

\begin{frame}[fragile]{Feature extraction and statistics}
\protect\phantomsection\label{feature-extraction-and-statistics-3}
\fontsize{9}{9}\sf

Find the most trended time series:

\begin{Shaded}
\begin{Highlighting}[]
\NormalTok{most\_trended }\OtherTok{\textless{}{-}}\NormalTok{ tourism }\SpecialCharTok{|\textgreater{}}
  \FunctionTok{features}\NormalTok{(Trips, feat\_stl) }\SpecialCharTok{|\textgreater{}}
  \FunctionTok{filter}\NormalTok{(trend\_strength }\SpecialCharTok{==} \FunctionTok{max}\NormalTok{(trend\_strength))}
\end{Highlighting}
\end{Shaded}

\pause\vspace*{-0.3cm}

\begin{Shaded}
\begin{Highlighting}[]
\NormalTok{tourism }\SpecialCharTok{|\textgreater{}}
  \FunctionTok{right\_join}\NormalTok{(most\_trended, }\AttributeTok{by =} \FunctionTok{c}\NormalTok{(}\StringTok{"State"}\NormalTok{, }\StringTok{"Region"}\NormalTok{, }\StringTok{"Purpose"}\NormalTok{)) }\SpecialCharTok{|\textgreater{}}
  \FunctionTok{ggplot}\NormalTok{(}\FunctionTok{aes}\NormalTok{(}\AttributeTok{x =}\NormalTok{ Quarter, }\AttributeTok{y =}\NormalTok{ Trips)) }\SpecialCharTok{+}
  \FunctionTok{geom\_line}\NormalTok{() }\SpecialCharTok{+} \FunctionTok{facet\_grid}\NormalTok{(}\FunctionTok{vars}\NormalTok{(State, Region, Purpose))}
\end{Highlighting}
\end{Shaded}

\pandocbounded{\includegraphics[keepaspectratio]{02_features_files/figure-beamer/extreme4-1.pdf}}
\end{frame}

\begin{frame}[fragile]{Feature extraction and statistics}
\protect\phantomsection\label{feature-extraction-and-statistics-4}
\fontsize{9}{10}\sf

\begin{Shaded}
\begin{Highlighting}[]
\NormalTok{tourism }\SpecialCharTok{|\textgreater{}} \FunctionTok{features}\NormalTok{(Trips, feat\_acf)}
\end{Highlighting}
\end{Shaded}

\begin{verbatim}
# A tibble: 304 x 10
   Region     State Purpose     acf1 acf10 diff1_acf1 diff1_acf10 diff2_acf1 diff2_acf10
   <chr>      <chr> <chr>      <dbl> <dbl>      <dbl>       <dbl>      <dbl>       <dbl>
 1 Adelaide   SA    Busine~  0.0333  0.131     -0.520       0.463     -0.676       0.741
 2 Adelaide   SA    Holiday  0.0456  0.372     -0.343       0.614     -0.487       0.558
 3 Adelaide   SA    Other    0.517   1.15      -0.409       0.383     -0.675       0.792
 4 Adelaide   SA    Visiti~  0.0684  0.294     -0.394       0.452     -0.518       0.447
 5 Adelaide ~ SA    Busine~  0.0709  0.134     -0.580       0.415     -0.750       0.746
 6 Adelaide ~ SA    Holiday  0.131   0.313     -0.536       0.500     -0.716       0.906
 7 Adelaide ~ SA    Other    0.261   0.330     -0.253       0.317     -0.457       0.392
 8 Adelaide ~ SA    Visiti~  0.139   0.117     -0.472       0.239     -0.626       0.408
 9 Alice Spr~ NT    Busine~  0.217   0.367     -0.500       0.381     -0.658       0.587
10 Alice Spr~ NT    Holiday -0.00660 2.11      -0.153       2.11      -0.274       1.55 
# i 294 more rows
# i 1 more variable: season_acf1 <dbl>
\end{verbatim}
\end{frame}

\section{Dimension reduction for
features}\label{dimension-reduction-for-features}

\begin{frame}[fragile]{Feature extraction and statistics}
\protect\phantomsection\label{feature-extraction-and-statistics-5}
\begin{Shaded}
\begin{Highlighting}[]
\NormalTok{tourism\_features }\OtherTok{\textless{}{-}}\NormalTok{ tourism }\SpecialCharTok{|\textgreater{}}
  \FunctionTok{features}\NormalTok{(Trips, }\FunctionTok{feature\_set}\NormalTok{(}\AttributeTok{pkgs =} \StringTok{"feasts"}\NormalTok{))}
\end{Highlighting}
\end{Shaded}

\fontsize{8}{8}\sf\vspace*{-0.3cm}

\begin{verbatim}
# A tibble: 304 x 51
   Region         State Purpose trend_strength seasonal_strength_year seasonal_peak_year
   <chr>          <chr> <chr>            <dbl>                  <dbl>              <dbl>
 1 Adelaide       SA    Busine~          0.464                  0.407                  3
 2 Adelaide       SA    Holiday          0.554                  0.619                  1
 3 Adelaide       SA    Other            0.746                  0.202                  2
 4 Adelaide       SA    Visiti~          0.435                  0.452                  1
 5 Adelaide Hills SA    Busine~          0.464                  0.179                  3
 6 Adelaide Hills SA    Holiday          0.528                  0.296                  2
 7 Adelaide Hills SA    Other            0.593                  0.404                  2
 8 Adelaide Hills SA    Visiti~          0.488                  0.254                  0
 9 Alice Springs  NT    Busine~          0.534                  0.251                  0
10 Alice Springs  NT    Holiday          0.381                  0.832                  3
# i 294 more rows
# i 45 more variables: seasonal_trough_year <dbl>, spikiness <dbl>, linearity <dbl>,
#   curvature <dbl>, stl_e_acf1 <dbl>, stl_e_acf10 <dbl>, acf1 <dbl>, acf10 <dbl>,
#   diff1_acf1 <dbl>, diff1_acf10 <dbl>, diff2_acf1 <dbl>, diff2_acf10 <dbl>,
#   season_acf1 <dbl>, pacf5 <dbl>, diff1_pacf5 <dbl>, diff2_pacf5 <dbl>,
#   season_pacf <dbl>, zero_run_mean <dbl>, nonzero_squared_cv <dbl>,
#   zero_start_prop <dbl>, zero_end_prop <dbl>, lambda_guerrero <dbl>, ...
\end{verbatim}

\begin{textblock}{4.5}(10.6,.5)
\begin{alertblock}{}\fontsize{10}{12}\sf
All features from the feasts package
\end{alertblock}
\end{textblock}
\end{frame}

\begin{frame}[fragile]{Feature extraction and statistics}
\protect\phantomsection\label{feature-extraction-and-statistics-6}
\begin{Shaded}
\begin{Highlighting}[]
\NormalTok{pcs }\OtherTok{\textless{}{-}}\NormalTok{ tourism\_features }\SpecialCharTok{|\textgreater{}}
  \FunctionTok{select}\NormalTok{(}\SpecialCharTok{{-}}\NormalTok{State, }\SpecialCharTok{{-}}\NormalTok{Region, }\SpecialCharTok{{-}}\NormalTok{Purpose) }\SpecialCharTok{|\textgreater{}}
  \FunctionTok{prcomp}\NormalTok{(}\AttributeTok{scale =} \ConstantTok{TRUE}\NormalTok{) }\SpecialCharTok{|\textgreater{}}
\NormalTok{  broom}\SpecialCharTok{::}\FunctionTok{augment}\NormalTok{(tourism\_features)}
\end{Highlighting}
\end{Shaded}

\fontsize{8}{8}\sf\vspace*{-0.3cm}

\begin{verbatim}
# A tibble: 304 x 100
   .rownames Region         State Purpose  trend_strength seasonal_strength_year
   <chr>     <chr>          <chr> <chr>             <dbl>                  <dbl>
 1 1         Adelaide       SA    Business          0.464                  0.407
 2 2         Adelaide       SA    Holiday           0.554                  0.619
 3 3         Adelaide       SA    Other             0.746                  0.202
 4 4         Adelaide       SA    Visiting          0.435                  0.452
 5 5         Adelaide Hills SA    Business          0.464                  0.179
 6 6         Adelaide Hills SA    Holiday           0.528                  0.296
 7 7         Adelaide Hills SA    Other             0.593                  0.404
 8 8         Adelaide Hills SA    Visiting          0.488                  0.254
 9 9         Alice Springs  NT    Business          0.534                  0.251
10 10        Alice Springs  NT    Holiday           0.381                  0.832
# i 294 more rows
# i 94 more variables: seasonal_peak_year <dbl>, seasonal_trough_year <dbl>,
#   spikiness <dbl>, linearity <dbl>, curvature <dbl>, stl_e_acf1 <dbl>,
#   stl_e_acf10 <dbl>, acf1 <dbl>, acf10 <dbl>, diff1_acf1 <dbl>, diff1_acf10 <dbl>,
#   diff2_acf1 <dbl>, diff2_acf10 <dbl>, season_acf1 <dbl>, pacf5 <dbl>,
#   diff1_pacf5 <dbl>, diff2_pacf5 <dbl>, season_pacf <dbl>, zero_run_mean <dbl>,
#   nonzero_squared_cv <dbl>, zero_start_prop <dbl>, zero_end_prop <dbl>, ...
\end{verbatim}

\begin{textblock}{4.5}(10.6,.5)
\begin{alertblock}{}\fontsize{10}{12}\sf
Principal components based on all features from the feasts package
\end{alertblock}
\end{textblock}
\end{frame}

\begin{frame}[fragile]{Feature extraction and statistics}
\protect\phantomsection\label{feature-extraction-and-statistics-7}
\fontsize{9}{9}\sf

\begin{textblock}{3.3}(1.4,3)
\begin{alertblock}{}\fontsize{10}{12}\sf
Principal components based on all features from the feasts package
\end{alertblock}
\end{textblock}

\begin{Shaded}
\begin{Highlighting}[]
\NormalTok{pcs }\SpecialCharTok{|\textgreater{}} \FunctionTok{ggplot}\NormalTok{(}\FunctionTok{aes}\NormalTok{(}\AttributeTok{x=}\NormalTok{.fittedPC1, }\AttributeTok{y=}\NormalTok{.fittedPC2)) }\SpecialCharTok{+}
  \FunctionTok{geom\_point}\NormalTok{() }\SpecialCharTok{+} \FunctionTok{theme}\NormalTok{(}\AttributeTok{aspect.ratio=}\DecValTok{1}\NormalTok{)}
\end{Highlighting}
\end{Shaded}

\placefig{7}{2.6}{height=6.4cm, width=12cm}{pca1}
\vspace*{10cm}
\end{frame}

\begin{frame}[fragile]{Feature extraction and statistics}
\protect\phantomsection\label{feature-extraction-and-statistics-8}
\fontsize{9}{9}\sf

\begin{textblock}{3.3}(1.4,3)
\begin{alertblock}{}\fontsize{10}{12}\sf
Principal components based on all features from the feasts package
\end{alertblock}
\end{textblock}

\begin{Shaded}
\begin{Highlighting}[]
\NormalTok{pcs }\SpecialCharTok{|\textgreater{}} \FunctionTok{ggplot}\NormalTok{(}\FunctionTok{aes}\NormalTok{(}\AttributeTok{x=}\NormalTok{.fittedPC1, }\AttributeTok{y=}\NormalTok{.fittedPC2, }\AttributeTok{col=}\NormalTok{State)) }\SpecialCharTok{+}
  \FunctionTok{geom\_point}\NormalTok{() }\SpecialCharTok{+} \FunctionTok{theme}\NormalTok{(}\AttributeTok{aspect.ratio=}\DecValTok{1}\NormalTok{)}
\end{Highlighting}
\end{Shaded}

\placefig{7}{2.6}{height=6.4cm, width=12cm}{pca2}
\vspace*{10cm}
\end{frame}

\begin{frame}[fragile]{Feature extraction and statistics}
\protect\phantomsection\label{feature-extraction-and-statistics-9}
\fontsize{9}{9}\sf

\begin{textblock}{3.3}(1.4,3)
\begin{alertblock}{}\fontsize{10}{12}\sf
Principal components based on all features from the feasts package
\end{alertblock}
\end{textblock}

\begin{Shaded}
\begin{Highlighting}[]
\NormalTok{pcs }\SpecialCharTok{|\textgreater{}} \FunctionTok{ggplot}\NormalTok{(}\FunctionTok{aes}\NormalTok{(}\AttributeTok{x=}\NormalTok{.fittedPC1, }\AttributeTok{y=}\NormalTok{.fittedPC2, }\AttributeTok{col=}\NormalTok{Purpose)) }\SpecialCharTok{+}
  \FunctionTok{geom\_point}\NormalTok{() }\SpecialCharTok{+} \FunctionTok{theme}\NormalTok{(}\AttributeTok{aspect.ratio=}\DecValTok{1}\NormalTok{)}
\end{Highlighting}
\end{Shaded}

\only<1>{\placefig{7}{2.6}{height=6.4cm, width=12cm}{pca3}}
\only<2>{\placefig{7}{2.6}{height=6.4cm, width=12cm}{pca4}}
\vspace*{10cm}
\end{frame}

\begin{frame}[fragile]{Feature extraction and statistics}
\protect\phantomsection\label{feature-extraction-and-statistics-10}
\fontsize{8}{8}\sf

\begin{Shaded}
\begin{Highlighting}[]
\NormalTok{outliers }\SpecialCharTok{|\textgreater{}}
  \FunctionTok{left\_join}\NormalTok{(tourism, }\AttributeTok{by =} \FunctionTok{c}\NormalTok{(}\StringTok{"State"}\NormalTok{, }\StringTok{"Region"}\NormalTok{, }\StringTok{"Purpose"}\NormalTok{)) }\SpecialCharTok{|\textgreater{}}
  \FunctionTok{mutate}\NormalTok{(}\AttributeTok{Series =} \FunctionTok{glue}\NormalTok{(}\StringTok{"\{State\}"}\NormalTok{, }\StringTok{"\{Region\}"}\NormalTok{, }\StringTok{"\{Purpose\}"}\NormalTok{, }\AttributeTok{.sep =} \StringTok{"}\SpecialCharTok{\textbackslash{}n\textbackslash{}n}\StringTok{"}\NormalTok{)) }\SpecialCharTok{|\textgreater{}}
  \FunctionTok{ggplot}\NormalTok{(}\FunctionTok{aes}\NormalTok{(}\AttributeTok{x =}\NormalTok{ Quarter, }\AttributeTok{y =}\NormalTok{ Trips)) }\SpecialCharTok{+}
  \FunctionTok{geom\_line}\NormalTok{() }\SpecialCharTok{+} \FunctionTok{facet\_grid}\NormalTok{(Series }\SpecialCharTok{\textasciitilde{}}\NormalTok{ .) }\SpecialCharTok{+}
  \FunctionTok{labs}\NormalTok{(}\AttributeTok{title =} \StringTok{"Outlying time series in PC space"}\NormalTok{)}
\end{Highlighting}
\end{Shaded}

\includegraphics[width=\linewidth,height=0.55\textheight,keepaspectratio]{02_features_files/figure-beamer/outliers2-1.pdf}
\end{frame}




\end{document}
